\section{Bancada de generación seleccionada y retirada segura (GCI-23)}
\label{sec:sd4-bancada-gci23}

Synaptix selecciona dos bancos Crestchic \emph{Datacentre 200kW Stackable Transportable}, uno por grupo, y un tercero de reserva transportable. Esta selección sustituye la Avtron 4100 del desarrollo precedente, incluida la atribución de ALC; VCS con LC10 es un control manual. La variante publicada opera a 400 V trifásicos y 50/60 Hz, con resolución de 1 kW, tolerancia de elementos de $\pm5\,\%$ y auxiliar independiente de 220--240 V monofásicos, 10 A, 50/60 Hz. Los 200 kW son capacidad de placa de la bancada: no son una carga autorizada para el P200-6. Se selecciona el modo de alimentación auxiliar separado, para conservar refrigeración tras retirar carga resistiva \parencite[pp.~1--2]{lote23-crest200}.

Cada unidad mide aproximadamente $1.183\times902\times983$ mm y pesa 275 kg; el aire horizontal publicado es 6.435 m$^3$/h y la elevación de temperatura de salida es menor de 100\,$^{\circ}$C. El recinto técnico de proyecto dispone cada bancada sobre base propia, con descarga al exterior sin recircular a la admisión del grupo, sin ocupar rutas de evacuación ni descargar sobre combustible o personas. El fabricante publica IP55 para la \emph{cámara de control}, no para toda la unidad; por ello se prevé alojamiento protegido de lluvia y salpicaduras. El dominio publicado es $-20$ a $+50\,^{\circ}$C y altitud menor de 500 m; las distancias de instalación, ruido combinado, corrosión y desempeño bajo humedad ambiental requieren el manual específico y su verificación de instalación. No se trasladan distancias de Avtron. Corte térmico, sobrecarga de ventilador, sobrecorriente y parada de emergencia son protecciones publicadas; no sustituyen el permiso externo por carga total del grupo \parencite[p.~1]{lote23-crest200}.

\textbf{Cota conjunta y aplicación manual supervisada.} Se conserva el límite propio de proyecto de 90 kVA/90 kW por grupo hasta que el balance integral autorice otra combinación. El auxiliar se reserva a 2,40 kVA por unidad activa, cota de $240(10)/1000$, sin convertirla en una potencia real de ventilador ni atribuirle un factor de potencia. A 400 V, la cota superior de un paso es 1,05 kW; para la variación de tensión breve de $+5\,\%$ publicada resulta $1,05(1,05)^2=1,157625$ kW. No se interpreta la tolerancia breve como autorización de operación continua a 420 V.

Para cada estado y tensión dentro del dominio, el permiso previo comprueba $S_{serv}+S_{aux}+S_{bancada}\leq90$ kVA, $P_{serv}+P_{aux}+P_{bancada}\leq90$ kW y corrientes de las tres fases; el valor solicitado al LC10 sólo se admite si incluye el paso siguiente completo y sus tolerancias. Con el cribado anterior de servicio de 74,824595 kVA, quedan $90-74,824595-2,40=12,775405$ kVA. Once pasos nominales tienen cota 12,733875 kW y el duodécimo eleva esa cota a 13,891500 kW, por lo que no se permite. Este máximo de once pasos conserva el cribado agregado anterior; GCI-26 reduce la solicitud vigente a diez pasos por el nuevo circuito de 16 A. Ninguna de esas cifras es consigna universal. Un receptor que se incorpora, la recarga o el solape climático pueden reducirlo a cero. No se ofrece regulación automática desde VCS ni se garantiza mantener 57,6 kW con esos once pasos. La baja carga admisible del motor continúa requiriendo criterio de FG Wilson y un perfil que evite acumular horas fuera de ese criterio.

La operación de ensayo mantiene el LC10 bajo control del operador autorizado y suspende toda adición ante transición de fuente, solicitud de transferencia, recarga, arranque climático, exceso de fase o pérdida de supervisión. La retirada por seguridad abre la alimentación resistiva antes de incorporar el servicio anunciado; el auxiliar se mantiene desde la fuente protegida durante el enfriamiento prescrito. La señal externa, medición de grupo, contactores, enclavamientos y tiempos hasta apertura completa deben quedar individualizados para liberar ese permiso. La lectura propia de VCS corresponde a la bancada y no demuestra por sí misma la corriente total del grupo. El tiempo de apertura de un accesorio no demuestra el tiempo total del sistema. El ensayo registrará escalón aplicado y retirado, tensión/frecuencia, máximos por fase, consumo de combustible y respuesta de los receptores; no se atribuyen resultados ejecutados.

\section{Retirada de rectificador por seguridad y condiciones de recuperación (CEN-23)}
\label{sec:sd4-retirada-cen23}

Para la retirada de emergencia se seleccionan cuatro disparadores de mínima tensión Schneider \texttt{LV429407}, 220--240 V AC/50--60 Hz, compatibles con la familia ComPacT NSX: uno en cada entrada AC de rectificador y uno en cada alimentación resistiva de bancada. Se conserva alimentación auxiliar separada para ventiladores. La bobina MN permanece energizada mientras exista permiso; parada de emergencia, alarma de hidrógeno, caudal insuficiente, fallo de circuito y pérdida de alimentación retiran permiso. Los caminos A y B tienen circuitos separados; no se usa una única orden que retire el camino sano sin evaluar su seguridad. La fuente de cada permiso se toma del camino protegido correspondiente, con protección propia y contactos de seguridad en serie, independientemente del PLC de registro. La alimentación y extracción de seguridad no pasan por el interruptor retirado del rectificador \parencite[pp.~C-37 y F-38]{lote23-nsx2026}.

En MN la apertura es garantizada con tensión $\leq0,35U_n$; entre $0,35$ y $0,70U_n$ no se garantiza, y el cierre requiere al menos $0,85U_n$. Una apertura franca del contacto de permiso evita diseñar el disparo mediante una caída parcial de tensión. El catálogo publica 50 ms de respuesta del accesorio, 10 VA de captación y 5 VA de mantenimiento; el capítulo de cableado aconseja aproximadamente 30 VA para captación. Se reserva esta última cota: cuatro bobinas representan 20 VA continuos y 120 VA de captación simultánea antes de relés y protecciones. La selección de contactos, fuente, caída de tensión, selectividad y tiempos de interruptor debe satisfacer ese circuito completo. Tras disparo se exige rearme manual con condiciones seguras restablecidas; no hay reconexión automática por volver la red. MN se reserva para seguridad, no para regular habitualmente carga o recarga: el fabricante advierte desgaste del mecanismo y reducción de su vida mecánica por uso repetido \parencite[pp.~C-37 y E-20]{lote23-nsx2026}.

La retirada AC del rectificador de 93T con bypass bloqueado elimina la fuente de recarga, pero consume batería mientras permanezca carga. El inversor no mantiene servicio indefinido después de esa maniobra. Antes de agotar la reserva se transfiere el servicio compatible al camino sano; las entradas únicas requieren la selección y sostenimiento propios de CEN-20. La extracción del banco afectado sigue habilitada y la generación de gas no se presume nula tras abrir AC. Este circuito selecciona un medio de seguridad externo; no atribuye a CN12 una inhibición selectiva de cargador no documentada ni sustituye el control de 54 A por cadena. La recepción debe verificar fallo abierto de cada permiso, extracción mantenida, rama sana, retirada resistiva, rearme y recarga seguida de nuevo corte \parencite{p7-eaton-93t-guia}.

\section{Demanda integral y frontera de selección (ECI-23)}
\label{sec:sd4-balance-eci23}

Los auxiliares de bancada se contabilizan durante ensayo o asistencia de generación, aguas arriba de UPS; sus 2,40 kVA por unidad activa no se añaden a la salida de cada 93T ni se duplican en el régimen normal. Las bobinas de seguridad, fuentes, detectores, tratamiento y conversión sí se incorporan a su alimentación efectiva, con potencia real y aparente diferenciadas. Los 20 VA de bobinas no equivalen a 20 W medidos. El reemplazo de bancada elimina la incompatibilidad de auxiliar de 60 Hz; no elimina el exceso de L1 calculado en ECI-20 ni dimensiona climatización por diferencia de potencia de placa.

La partida parcial de 50,360015 kW/54,432488 kVA, con margen de compra, deja 9,639985 kW/5,567512 kVA respecto de cada 93T60 antes de tratamiento. Si se integran dos recalentadores de 9 kW como carga simultánea de proyecto, su sola potencia llevaría esa partida a 68,360015 kW sin margen adicional sobre calentadores: supera 60 kW. Si se aplica también el margen del veinte por ciento a esos 18 kW, alcanza 71,960015 kW antes de controles. ENV-25 añade 0,100 kW de reserva de TTC25, con veinte por ciento 0,120 kW: la cota activa asciende a 72,080015 kW antes del resto de tratamiento. Los 76,0485 kVA previos no incluyen estos controles; su potencia aparente y régimen de triac deben incorporarse por estado, sin asumir factor de potencia unitario del conjunto. Por consiguiente, no se libera una selección de 93T60 para todo el parque mediante un ajuste de fase o recarga; el frío, bombas y unidades de aire sólo aumentan esa frontera. Esas cifras son cribado de combinación, no consumo operacional ni datos de catálogo de una variante climática.

Una sustitución de UPS exige resolver simultáneamente potencia de salida por fase, recarga documentada, autonomía del parque completo, número de cadenas, implantación, ventilación, extinción, entrada de generador y protecciones. El balance por estados separa red, batería, grupo, retorno/recarga y ensayo y aplica las pérdidas de la variante finalmente integrada. El perfil anual utiliza horas excluyentes y consumo real de auxiliares por estado; el factor de potencia y el PUE no se deducen de cocientes de cotas de compra ni de sumar dos caminos supervivientes. Los litros de combustible no se suman al numerador eléctrico. El informe de diseño final incluirá estas fronteras y la recepción medirá por separado TI, clima, conversión y exportación; los valores previos de PUE no constituyen un resultado actualizado del recinto.
